\usepackage[T1]{fontenc}

% This assumes your files are encoded as UTF8
\usepackage[utf8]{inputenc}

% This is not strictly necessary, and may be commented out,
% but it will improve the layout of the manuscript,
% and will typically save some space.
\usepackage{microtype}

% This is also not strictly necessary, and may be commented out.
% However, it will improve the aesthetics of text in
% the typewriter font.
\usepackage{inconsolata}

%Including images in your LaTeX document requires adding
%additional package(s)
\usepackage{graphicx}

% Recommended, but optional, packages for figures and better typesetting:
\usepackage{graphbox} %loads graphicx package
\usepackage{subcaption} % for subfigures
\usepackage{paralist}
% Load xcolor with the table option before colortbl for proper \rowcolors support
\usepackage[table]{xcolor}
\usepackage{booktabs} % for professional tables
\usepackage{colortbl} % For coloring rows in tables
\usepackage{makecell} % For customizing table cells
\usepackage{multirow} % For multirow cells in tables
% \usepackage{color} % removed: xcolor (loaded above with [table]) supersedes color
\usepackage{amsmath}
\usepackage[normalem]{ulem} % for strikethrough

\usepackage{longtable}

% hyperref makes hyperlinks in the resulting PDF.
% If your build breaks (sometimes temporarily if a hyperlink spans a page)
% please comment out the following usepackage line and replace
\usepackage{hyperref}

% Algorithm packages
\usepackage{algorithm}
\usepackage{algorithmic}

% Attempt to make hyperref and algorithmic work together better:
% (algorithm package already defines \theHalgorithm, so we skip the manual definition)

% NeurIPS 2025 Conference style
\usepackage[preprint]{neurips_2025}

% For theorems and such
\usepackage{amssymb}
\usepackage{mathtools}
\usepackage{amsthm}

\usepackage{verbatimbox}
\usepackage{fvextra}

% if you use cleveref..
\usepackage[capitalize,noabbrev]{cleveref}

% Add alltt package for verbatim text with special characters
\usepackage{alltt}

%%%%%%%%%%%%%%%%%%%%%%%%%%%%%%%%
% THEOREMS
%%%%%%%%%%%%%%%%%%%%%%%%%%%%%%%%
\theoremstyle{plain}
\newtheorem{theorem}{Theorem}[section]
\newtheorem{proposition}[theorem]{Proposition}
\newtheorem{lemma}[theorem]{Lemma}
\newtheorem{corollary}[theorem]{Corollary}
\theoremstyle{definition}
\newtheorem{definition}[theorem]{Definition}
\newtheorem{assumption}[theorem]{Assumption}
\theoremstyle{remark}
\newtheorem{remark}[theorem]{Remark}

% Todonotes is useful during development; simply uncomment the next line
%    and comment out the line below the next line to turn off comments
\usepackage[disable,textsize=tiny]{todonotes}
% \usepackage[textsize=tiny]{todonotes}

% Boxes for results
\usepackage{boxedminipage}

% Add packages for code listings
\usepackage{listings}
\lstset{
    breaklines=true,
    basicstyle=\ttfamily,
    keepspaces=true,
}

\usepackage{xspace}
% \usepackage{longtable} % removed: duplicate of line 34

% TikZ for pipeline diagrams
\usepackage{tikz}
\usetikzlibrary{shapes.geometric, arrows.meta, positioning, fit, backgrounds, calc, shadows.blur}

% Custom environment for prompts
\usepackage[breakable,skins]{tcolorbox}
\usepackage{adjustbox}

% Define a custom environment for prompts
\newtcolorbox{promptbox}[1][]{
    colback=gray!10,
    colframe=gray!20,
    boxrule=0pt,
    arc=0pt,
    outer arc=0pt,
    left=2pt,
    right=2pt,
    top=2pt,
    bottom=2pt,
    boxsep=2pt,
    % Change to a smaller font:
    before upper={\scriptsize\ttfamily},
    breakable,
    #1
}

% We will caption outside the float environment:
\usepackage{capt-of}  % for \captionof{figure}
\usepackage{float}

% Configure fancyvrb default settings
\usepackage{fancyvrb}
\fvset{
  frame=none,
  numbers=left,
  numbersep=3pt,
  formatcom=\color{black}
}

% Disable special character handling in verbatim
\makeatletter
\def\FV@ProcessLine#1{\hbox to \linewidth{\hskip\@totalleftmargin\FV@LeftListNumber\FV@LeftListFrame \FancyVerbFormatLine{#1}\hfil\FV@RightListFrame\FV@RightListNumber}}
\makeatother

% Configure listings for better line number display with preserved empty lines
\lstdefinestyle{numberedcode}{
  basicstyle=\ttfamily\scriptsize,
  numbers=left,
  numberstyle=\tiny\color{black},
  stepnumber=1,
  numbersep=5pt,
  backgroundcolor=\color{gray!10},
  frame=none,
  breaklines=true,
  showstringspaces=false,
  columns=fullflexible,
  keepspaces=true,
  emptylines=1,
  xleftmargin=15pt,
  escapechar=@
}

% Completely redesigned promptverbatim environment
\newenvironment{promptverbatim}{%
  \begin{figure}[t]
  \begin{minipage}{\linewidth}
}{%
  \end{minipage}
  \end{figure}
}

\newcommand{\captionprompt}[1]{%
  \caption{#1}
  \label{prompt:#1}
}

% New command to include code with line numbers
\newcommand{\codewithlinenumbers}[1]{%
  \begin{lstlisting}[style=numberedcode]
#1
  \end{lstlisting}
}